\section{Základný popis aplikácie}

\subsection{Motivácia}
V poslednej dobe narastá počet investicií, ktoré smerujú do nehnuteľností a nie vždy sú využívané širokou verejnosťou, ale iba úzkou skupinou ľudí, napríklad rodinami, priateľmi alebo komunitami.
Najväčším problémom rezervačných systémov ako booking.com, airbnb.com, či vrbo.com je ich zameranie na širokú verejnosť s vysokými poplatkami, nedostatočnou kontrolou nad rezerváciami a obmedzenými možnosťami prispôsobenia pre majiteľov nehnuteľností. 
Navyše, tieto platformy často neberú do úvahy špecifické potreby a požiadavky majiteľov nehnuteľností, čo môže viesť k frustrácii a neefektívnemu spravovaniu rezervácií.

\subsection{Ciele našej aplikácie}
Cieľom našej aplikácie je vytvoriť privátny rezervačný systém, ktorý bude slúžiť pre konkrétnu skupinu užívateľov, napríklad pre rodinu, priateľov, malú komunitu alebo aj širšiu verejnosť.
Hlavnou motiváciou je poskytnúť majiteľovi prehľad o obsadenosti apartmánov a nehnuteľností, umožniť jednoduché spravovanie rezervácií a odstrániť závislosť na externých platformách, ktoré môžu byť nákladné a neflexibilné.

\subsection{Cieľová skupina}
Našou cieľovou skupinou sú majitelia nehnuteľností, ktorí hľadajú efektívny a flexibilný spôsob spravovania rezervácií bez závislosti na externých platformách. 
Títo majitelia môžu byť jednotlivci, rodiny, priatelia alebo malá komunita, ktorá chce mať kontrolu nad svojimi rezerváciami a nechcú platiť vysoké poplatky externým platformám.

\subsection{Desktop vs Mobile}
Dizajn aplikácie bude responzívny tak, aby bola aplikácia prístupná a plnohodnotne použiteľná na desktopových zariadeniach, tabletoch aj mobilných telefónoch.
Na desktopových zariadeniach bude hlavné menu zobrazené horizontálne, zatiaľ čo na mobilných zariadeniach bude riešené vo vertikálnej podobe.
Obsah jednotlivých stránok bude usporiadaný pomocou grid layoutu, ktorý umožní flexibilné rozloženie prvkov podľa veľkosti obrazovky a zabezpečí dobrú použiteľnosť na všetkých podporovaných zariadeniach.
Na menších obrazovkách budú jednotlivé sekcie a ovládacie prvky usporiadané pod sebou tak, aby používateľ vedel aplikáciu pohodlne ovládať aj dotykovo.